\chapter{Stato, che cosa resta, e riproduzione}
\label{ch:status}

Questo capitolo è il registro. Elenca ogni suite di test che il
programma possiede e il suo verdetto, dichiara esattamente quali
grandezze un'esecuzione di progetto può fissare e quali può muovere,
dice che cosa non è ancora stato costruito e fornisce il comando che
riesegue ogni risultato da zero. Nulla qui è fisica nuova: è la parte
che serve a un lettore per verificare il resto.

\section{Stato di ogni suite}

\begin{table}[htbp]
\centering
\small
\begin{tabular}{llll}
\toprule
Passo & Carrier (\code{validation/}) & Verdetto & Capitolo \\
\midrule
Marcia del Mattone 1 & \code{a1\_ideal\_march\_jax.py} & certificata &
  (prereq.) \\
Funzionale di spinta & \code{a1\_thrust\_functional.py} & 8/8 PASS &
  \ref{ch:score} \\
Strumento di spigolo & \code{a1\_corner\_instrument.py} & 5/5 PASS &
  \ref{ch:score} \\
Driver in $\varepsilon$ vs forma chiusa & \code{a1\_driver\_eps.py} &
  4/4 PASS & \ref{ch:firstopt} \\
Marcia a parete assegnata & \code{a1\_wall\_march.py} & 5/5 + 2/2 PASS &
  \ref{ch:wallinput} \\
Verifica fra codici GENO (O3.4) & \code{g0\_geno\_crosscode.py} & PASS &
  \ref{ch:referee} \\
Moltiplicatori di Rao vs AD (O3.3) & \code{a1\_o33\_toc.py} & 9/9 PASS &
  \ref{ch:referee} \\
Strato di ciclo (collasso + RDE simulato) & \code{a1\_cycle\_layer.py} &
  5/5 + 6/6 PASS & \ref{ch:cycle} \\
Misuratore di spigolo di ciclo (voti) & \code{a1\_cycle\_corner.py} & 8/8 PASS &
  \ref{ch:cycle} \\
Cella di getto libero & \code{a1\_freejet\_unit.py} & 6/6 PASS &
  \ref{ch:plug} \\
Marcia del plug (oracolo + assi) & \code{a1\_plug\_march.py} & 6/6 + 5/5 PASS &
  \ref{ch:plug} \\
Plug troncato sotto ciclo & \code{a1\_plug\_cycle.py} & 5/5 + 9/9 PASS &
  \ref{ch:plugcycle} \\
Gemello GENO a campo intero & \code{a1\_geno\_plug\_twin.py} & 8/8 PASS &
  \ref{ch:genotwin} \\
Marcia con swirl a vortice libero & \code{a1\_swirl\_march.py} & 9/9 PASS &
  \ref{ch:swirl} \\
Ingresso generale (rotazionale) & \code{a1\_rot\_march.py} & 8/8 PASS &
  \ref{ch:generalinlet} \\
Oracolo di flusso radiale ($v<0$) & \code{a1\_source\_flow\_oracle.py} & 6/6 PASS &
  \ref{ch:compare} \\
Confronto fra configurazioni & \code{a1\_config\_compare.py} & 7/10
  (3 aperte) & \ref{ch:compare} \\
Angolo di ingresso come variabile & \code{a1\_inlet\_angle\_opt.py} & 6/6 PASS &
  \ref{ch:compare} \\
Spina a forma libera (spline + TR-SQP) & \code{a1\_plug\_spline\_opt.py} &
  7/8 (vedi sotto) & \ref{ch:spline} \\
Nodi adattivi sulla spina & \code{a1\_plug\_adaptive.py} &
  11/12 (vedi sotto) & \ref{sec:adaptive} \\
Risoluzione del guadagno (scala appaiata) & \code{a1\_plug\_gain\_resolve.py} &
  RISOLTO: guadagno $= 0$ in banda & \ref{sec:gainresolved} \\
Misuratore di campo, due vie (CFD P0) & \code{a1\_field\_meter.py} &
  3/4 (1 aperta) & \ref{sec:fmeter} \\
\bottomrule
\end{tabular}
\caption{Ogni suite certificata del Mattone~2, il suo carrier e il suo
verdetto di riferimento. ``$m/n$ PASS'' conta le verifiche della
suite, rigettatori inclusi. Due suite portano verifiche aperte anziché
assorbirle in bande allargate: le chiusure di massa e quantità di moto
del confronto fra configurazioni per il plug discendente, e l'accordo
fra le due vie del misuratore di campo. Il guadagno un tempo aperto
della spina a forma libera (il suo 7/8 e l'11/12 del carrier adattivo
registrano i verdetti delle rispettive sessioni) è ora RISOLTO dal
carrier a scala appaiata: zero entro una banda dello $0.017\%$
(\S\ref{sec:gainresolved}). Le prime due sono la stessa fisica, e la
\S\ref{sec:fmeter} affina la diagnosi precedente: la vecchia cifra era
gonfiata dal misuratore, ma un residuo vicino all'$1\%$ sopravvive su
una superficie adatta a essere integrata e sopravvive al raffinamento,
quindi è una proprietà del campo e non solo della contabilità.
\code{a1\_flux\_meter.py} non porta un verdetto proprio --- è il modulo
strumento su cui quelle verifiche vanno ricablate, non una suite.}
\label{tab:status}
\end{table}

\section{Ingressi di progetto e vincoli, a oggi}
\label{sec:designinputs}

\subsection{Che cosa un'esecuzione di progetto può fissare}

Progettare significa tenere fisse alcune grandezze e lasciare che
l'ottimizzatore muova le altre. Ogni voce qui sotto può essere tenuta
fissa in modo indipendente per un'esecuzione; ciò che non è tenuto
fisso è ciò che l'ottimizzazione è autorizzata a cambiare.

\emph{Gas e camera}: le tabelle di proprietà (qualunque modello di gas
entro il suo intervallo di validità), le condizioni di camera
$(P_0, T_0)$ --- la pressione e la temperatura che il gas avrebbe a
riposo --- e, per le esecuzioni RDE, l'intero ciclo
$P_0(\xi), T_0(\xi)$ con il peso assegnato a ciascuna fase $\xi$ del
ciclo.

\emph{Ambiente}: la pressione ambiente $p_a$ in cui l'ugello scarica,
vuoto compreso. Per un ugello a spina non è mera contabilità: entra
nel flusso stesso, poiché la frontiera di getto libero è la superficie
su cui $p = p_a$.

\emph{Geometria di gola}: il raggio di gola $y_t$ e i due raggi di
raccordo $r_{tu}, r_{td}$ che arrotondano la gola, tutte manopole
differenziabili. Tenere fisso $y_t$ \emph{è} il vincolo di portata:
una gola bloccata lascia passare una portata fissa $\dot m$, quindi
fissarne il raggio fissa la portata del motore.

\emph{Dimensione complessiva}: il rapporto di espansione $\varepsilon$
--- l'area di uscita divisa per l'area di gola, l'unico numero che
dice quanto l'ugello espande il gas (il driver a una manopola dei
Capitoli~\ref{ch:firstopt} e~\ref{ch:cycle}).

\emph{La parete}: un'uscita quando il motore genera un contorno; un
ingresso interamente assegnato nelle modalità inverse a campana e a
plug; e un'incognita parametrizzata quando viene ottimizzata --- oggi
piccole protuberanze con gli estremi bloccati, punti di controllo
spline in seguito. Bloccare l'estremo \emph{è} il vincolo di lunghezza
fissa: è esattamente il problema che Rao pose, ed esattamente ciò che
il test di stazionarietà della Sezione~\ref{sec:o33} teneva fisso.

\emph{Modalità}: piana o assialsimmetrica; un singolo punto di
funzionamento stazionario o una media di ciclo. \emph{Risoluzione}: il
numero di raggi del ventaglio, di stazioni e di righe iniziali --- non
variabili di progetto, ma fissate per ogni esecuzione e parte di ogni
banda d'errore derivata.

\subsection{Quali vincoli si possono tenere simultaneamente}

Sul lato JAX i vincoli sono oggi imposti \emph{per costruzione}
anziché da un ottimizzatore vincolato generale: la portata tenendo la
gola, la lunghezza bloccando il labbro, l'ambiente come dato posto ---
e l'ottimizzatore muove ciò che resta (oggi $\varepsilon$, o la parete
attraverso la sua base di protuberanze). Tenere più vincoli non
lineari insieme contro il gradiente AD richiede la macchina generale:
un driver SQP a regione di fiducia --- un ottimizzatore vincolato
standard che a ogni iterazione compie un passo di dimensione limitata
e impone i vincoli al primo ordine --- che è nella lista del lavoro
restante, e per il quale sono stati costruiti i percorsi certificati a
parete assegnata.

Sul lato GENO il modulo RaoPlug (il suo caso di tipo 8) accetta, dalla
sessione a due vincoli, esattamente \emph{due} vincoli simultanei. Il
primo è il \emph{vincolo di massa}, sempre attivo: la curva di
progetto è terminata dove la massa accumulata al suo interno eguaglia
il bersaglio della gola $p_c\,\pi y_t^2/c^{*}$, e quella condizione è
consumata dal punto $D$ in cui la spina viene tagliata. Il secondo è
\emph{un} vincolo di progetto a scelta dell'utente, consumato da una
bisezione sull'angolo di uscita $\theta_E$ --- rapporto di espansione,
lunghezza del plug, angolo di uscita, spinta, numero di Mach di
uscita, pressione ambiente al labbro o pressione di base $p_b$
(confrontata con il valore che la condizione di spigolo implica al
taglio fissato dalla massa, l'accoppiamento sensato con il vincolo di
massa).

Due non è un limite implementativo ma \emph{teorico}. La famiglia di
spine di Rao ha esattamente due gradi di libertà a raggio di labbro
fissato --- l'angolo di uscita $\theta_E$ e il punto di taglio $D$ ---
quindi due condizioni la esauriscono, e una terza sarebbe
sovradeterminata. L'unica via onesta a un terzo vincolo simultaneo è
liberare il raggio di labbro $y_E$ come grado di libertà di scala, una
bisezione esterna già abbozzata nelle note di progetto di GENO e non
ancora costruita.

\section{Che cosa resta, in ordine}

Il filone GENO è ora chiuso da capo a fondo. La decisione del
proprietario era di risolvere il problema a due vincoli, e la sessione
GENO l'ha centrata: la curva di progetto termina alla massa di gola
assegnata, il punto di taglio fornisce il secondo grado di libertà con
la condizione di spigolo che restituisce la pressione di base che ne
consegue, quattro difetti di rigore sono stati corretti lungo la
strada e tre verifiche indipendenti PASSANO (la dualità fra i due modi
di scrivere la condizione terminale, $2\times10^{-5}$; il limite piano
esatto, bit per bit; un'integrazione Python indipendente della curva
di progetto, tre-quattro cifre). Una voce non si è chiusa: i numeri
tabulati nell'articolo del 1961 sulla spina, così come il registro li
cita, si sono rivelati non verificabili e internamente incoerenti (la
N-74 di GENO), quindi il proprietario deve procurare l'articolo
originale. Il confronto a campo intero --- la nostra marcia del plug
eseguita sulla spina propria di GENO --- è ora certificato 8/8 nel
Capitolo~\ref{ch:genotwin}, con la differenza misurata di
tre-quattro cifre fra i due codici caratterizzata e la sua causa
depositata come questione aperta dal lato GENO
(Sezione~\ref{sec:twingap}).

L'estensione allo swirl a vortice libero è certificata a livello di
marcia nel Capitolo~\ref{ch:swirl}: la garanzia strutturale misurata
anziché assunta, una soluzione esatta di condotto riprodotta e la
regressione a swirl spento a $10^{-13}$. Usarla dentro gli studi di
plug e di ciclo è ormai una scelta di posizione del problema, non
macchina nuova. Oltre il vortice libero il flusso è rotazionale ---
e il Capitolo~\ref{ch:generalinlet} porta ora quel caso dentro la
marcia, con gli invarianti di linea di corrente trasportati e
l'ingresso $(p, T, M, \theta)$ esattamente determinato.

La parete a forma libera, elencata qui per parecchie sessioni come
l'ultimo pezzo mancante, è ora costruita: il Capitolo~\ref{ch:spline}
dà alla spina una base di spline cubiche, un gradiente aggiunto e un
driver a regione di fiducia, e il risultato è un contorno che il
programma calcola anziché ereditare. Da come è stato costruito
seguono due cose. Usa direttamente le routine spline
dell'ottimizzatore a \emph{campana} anziché una base propria, quindi
campana e plug portano ora la stessa libertà di forma e il confronto
del Capitolo~\ref{ch:compare} diventa un confronto a pari condizioni.
E la scala di rappresentazione della Sezione~\ref{sec:reprladder}
mostra l'errore del plug concentrarsi nell'unica regione in cui il suo
contorno si piega, il che è un argomento a favore del posizionamento
adattivo dei nodi anziché di un maggior numero di nodi.

\begin{enumerate}[itemsep=2pt]
\item \textbf{Risolvere il guadagno della spina a forma libera} ---
  RISOLTO (\S\ref{sec:gainresolved}): su una scala appaiata a sei
  pioli fino a $(1921,1601)$, con l'incumbente riottimizzato a
  $(121,101)$ e il modello costretto a prevedere l'ultimo piolo prima
  di vederlo, il limite del guadagno è
  $-0.011\% \pm 0.017\%$ --- zero entro la banda. La linea di corrente
  troncata è ottimale a questa lunghezza e a questa ripartizione di
  ingresso entro lo $0.028\%$; ogni guadagno misurato in precedenza
  era strumento.
\item \textbf{Posizionamento adattivo dei nodi} sulla spina --- FATTO
  (\S\ref{sec:adaptive}); la sua verifica di verdetto è la voce
  precedente, ora chiusa.
\item \textbf{Rieseguire il confronto fra configurazioni} con
  entrambi i lati a forma libera, ora che il plug ha una base --- e
  solo una volta sistemata la voce precedente, poiché un margine
  citato contro un guadagno irrisolto ne erediterebbe il problema.
\item \textbf{Il confronto con Rao a piena espansione}: porre la
  costruzione classica e questo ottimizzatore alla STESSA lunghezza
  ($5.825$ m) con massa E ambiente adattati (\S\ref{sec:vsrao}). A
  $2.5$ m i due codici descrivono macchine diverse e le loro spinte
  non sono confrontabili.
\item \textbf{Il residuo di quantità di moto dell'$1\%$ del plug}
  (\S\ref{sec:fmeter}): ora misurato su una buona superficie e
  sopravvissuto al raffinamento, quindi è la questione fisica in
  sospeso di questo ramo.
\item \textbf{La superficie di controllo nel banco O3.3}: la nostra
  caratteristica terminale è presa come un'intera colonna di griglia,
  selezionata minimizzando proprio la dispersione che la verifica poi
  riporta. La dispersione misurata ($4.0\times10^{-3}$) non mostra
  patologie, ma il luogo va fissato \emph{a priori} e fermato al
  confine del nucleo anziché scelto dalla statistica che alimenta.
\item \textbf{Ingresso da ciclo misurato}: uno scambio di dati ---
  attenzione al rigettatore di intervallo di tabella del
  Capitolo~\ref{ch:cycle}.
\end{enumerate}

\section{Su quale solutore poggiano questi numeri}
\label{sec:rebase}

Ogni risultato di questo documento è stato prodotto originariamente su
una versione del solutore di flusso sottostante, ed è stato poi
rieseguito su un'altra. L'anello di Newton del solutore è stato
cambiato in modo indipendente: dove un tempo compiva un numero fisso
di iterazioni per cella, ora si ferma quando il passo scende sotto il
limite di certificazione della cella stessa. È un cambiamento reale al
modo in cui l'aritmetica termina, e sta sotto ogni numero qui, quindi
non è stato assunto innocuo.

È stato invece misurato. Rieseguire su entrambe le versioni il caso
che fissa la base comune dà una portata concorde \emph{bit per bit}, e
spinta, Mach di uscita, velocità di uscita, pressione di uscita e
raggio di labbro concordi fra $10^{-15}$ e $10^{-13}$ relativi ---
precisione di macchina, su quattro rapporti di espansione.
L'interpretazione è immediata: l'anello a iterazioni fisse stava già
convergendo oltre il punto in cui la nuova regola di arresto si ferma,
quindi i due concordano fino all'ultima cifra che ciascuno può
portare.

Nessun numero di questo documento si muove di conseguenza. Ogni suite
elencata nella Tabella~\ref{tab:status} è stata rieseguita sul nuovo
solutore e ha restituito il proprio verdetto registrato.

\section{Riproduzione}

Ogni carrier è autosufficiente e stampa il proprio verdetto; le
esecuzioni sono memorizzate dove indicato, quindi le riesecuzioni
costano poco. Sull'host \code{s2}:

\begin{lstlisting}
cd RDE_Design-rde-nozzle-program
.venv-a1/bin/python validation/a1_ideal_march_jax.py    # Mattone 1
.venv-a1/bin/python validation/a1_thrust_functional.py  # punteggio e bussola
.venv-a1/bin/python validation/a1_corner_instrument.py  # punteggio e bussola
.venv-a1/bin/python validation/a1_driver_eps.py         # prima ottimizzazione
.venv-a1/bin/python validation/a1_wall_march.py         # la parete come ingresso
.venv-a1/bin/python validation/g0_geno_crosscode.py     # l'arbitro
.venv-a1/bin/python validation/a1_o33_toc.py            # l'arbitro
.venv-a1/bin/python validation/a1_cycle_layer.py        # lo strato di ciclo
.venv-a1/bin/python validation/a1_cycle_corner.py       # lo strato di ciclo
.venv-a1/bin/python validation/a1_freejet_unit.py       # marcia del plug (cella)
.venv-a1/bin/python validation/a1_plug_march.py         # marcia del plug
.venv-a1/bin/python validation/a1_plug_cycle.py         # plug troncato + ciclo
.venv-a1/bin/python validation/a1_geno_plug_twin.py     # gemello plug GENO
.venv-a1/bin/python validation/a1_swirl_march.py        # swirl a vortice libero
.venv-a1/bin/python validation/a1_rot_march.py          # ingresso generale
.venv-a1/bin/python validation/a1_source_flow_oracle.py # oracolo radiale
.venv-a1/bin/python validation/a1_config_compare.py     # configurazioni
.venv-a1/bin/python validation/a1_inlet_angle_opt.py    # angolo di ingresso
.venv-a1/bin/python validation/a1_plug_spline_opt.py    # spina a forma libera
.venv-a1/bin/python validation/a1_plug_adaptive.py      # nodi adattivi
.venv-a1/bin/python validation/a1_plug_gain_resolve.py  # risoluzione del guadagno
.venv-a1/bin/python validation/a1_field_meter.py        # misuratore di campo (CFD P0)
\end{lstlisting}

Due di questi hanno bisogno di dati che gli script non generano da
sé. Il gemello del plug GENO legge un campo memorizzato dal codice
Fortran indipendente (\code{validation/\_geno\_twin/}), e il banco
O3.3 legge un contorno progettato dalla stessa fonte. Entrambi sono
prodotti di un codice che questo documento non modifica, quindi sono
ingressi anziché risultati. Il banco O3.3 porta anche un secondo
stadio che l'invocazione predefinita non esegue:

\begin{lstlisting}
STAGE=stationarity .venv-a1/bin/python validation/a1_o33_toc.py
\end{lstlisting}

Questo documento stesso è costruito da
\code{docs/brick2\_doc\_src/}: \code{make\_figures.py} rigenera le
figure di dati a partire dai dati di esecuzione memorizzati, e
\code{build.sh} compila \code{main.tex} (motore Tectonic) e installa
il PDF come \code{docs/rde\_nozzle\_A1\_brick2\_thrust.pdf}.

\section{Il misuratore di campo, e il primo cancello del programma
CFD}
\label{sec:fmeter}

\code{a1\_field\_meter.py} calcola la spinta per due vie indipendenti
su qualunque campo assialsimmetrico: come integrale di pressione
sull'hardware (il flusso di quantità di moto all'ingresso più
l'integrale di parete) e come flusso di quantità di moto attraverso un
taglio verticale a valle. Non condividono né superficie, né quadratura,
né integrando, quindi la loro differenza è una misura di conservazione
anziché di contabilità.

Esiste per due ragioni che si rivelano una sola. Le verifiche non
chiuse del confronto fra configurazioni avevano bisogno di un
misuratore adatto a essere integrato. E la controparte CFD di questo
programma (\code{docs/rde\_nozzle\_CFD\_plan.md}) ha bisogno del
proprio primo cancello: un obiettivo valutabile su un campo CFD e su
un campo di marcia che sia lo \emph{stesso funzionale}, in un solo
gauge e con un solo stato di riferimento. Confrontare una spinta CFD
con una spinta di marcia calcolata diversamente riprodurrebbe, a scala
molto maggiore, l'errore che questo documento ha ormai incontrato tre
volte.

Verdetto $3/4$. La via del flusso è indipendente dalla superficie
attraverso cui è presa ($0.40\%$ su quattro stazioni); il misuratore
riproduce la spinta registrata per lo stesso caso a
$1.0\times10^{-4}$; e un rigettatore conferma che due vie calcolate in
\emph{gauge ambiente diversi} non concordano, separando di
$5.8\times$.

La verifica che non passa è l'accordo fra le due vie stesso:
$0.99\%$, contro una banda dello $0.09\%$ che il raffinamento della
marcia sostiene, e la differenza si muove appena sotto raffinamento
($0.99\% \to 0.97\%$). È dunque una proprietà del campo. Questo affina
anziché contraddire il Capitolo~\ref{ch:compare}: l'$1.49\%$ di quel
capitolo era misurato sull'ultima colonna della marcia ed era gonfiato
da essa, ma circa l'$1\%$ è reale e resta la questione fisica in
sospeso del ramo del plug.
